% ch4_answers.tex -- answers to the Quick check questions in Module 6, Chapter 4.
\begin{chapteranswers}{4}{Transformers and Power Transmission}
  % ---- 4.1 The ideal transformer
  \item $N_\mathrm{s} = 2000$ is more than $N_\mathrm{p} = 400$, so it is a
        \final{step-up} transformer. Rearrange
        $\dfrac{V_\mathrm{p}}{V_\mathrm{s}} = \dfrac{N_\mathrm{p}}{N_\mathrm{s}}$:
        $V_\mathrm{s} = V_\mathrm{p}\dfrac{N_\mathrm{s}}{N_\mathrm{p}}
        = \qty{24}{V} \times \dfrac{2000}{400} = \qty{120}{V}$.
        Answer: \final{$V_\mathrm{s} = \qty{120}{V}$}.
  \item The peak scales by the turns ratio:
        $V_\mathrm{s} = \qty{9.0}{V} \times \dfrac{100}{250} = \qty{3.6}{V}$.
        A transformer does not change the frequency.
        Answer: \final{peak \qty{3.6}{V}} and \final{\qty{50}{Hz}}.
  \item (a) The secondary emf follows the \emph{gradient} of the current graph
        ($\varepsilon = -N\dfrac{\Delta\Phi}{\Delta t}$, and the flux follows the current).
        From 0 to \qty{5}{ms}, the current rises at a steady rate, so the emf is
        constant. From \qty{5}{ms} to \qty{10}{ms}, the current is steady, so the emf is
        zero. From \qty{10}{ms} to \qty{15}{ms}, the current falls at the same steady
        rate, so the emf is constant again: the same size, but the opposite sign.
        \par\smallskip{\centering
        \begin{tikzpicture}
          \begin{axis}[hsc, name=chFourAtop, width=72mm, height=28mm, xmin=0, xmax=20,
              ymin=-0.3, ymax=2.6, xtick={0,5,10,15,20}, xticklabels={}, ytick={0,2},
              ylabel={$I_\mathrm{p}$ (\si{A})}]
            \addplot[hscseriesA, line width=1.2pt] coordinates
              {(0,0) (5,2) (10,2) (15,0) (20,0)};
          \end{axis}
          \begin{axis}[hsc, at={(chFourAtop.south)}, anchor=north, yshift=-3mm,
              width=72mm, height=28mm, xmin=0, xmax=20, ymin=-1.5, ymax=1.5,
              xtick={0,5,10,15,20}, ytick=\empty, xlabel={Time $t$ (\si{ms})},
              ylabel={$\varepsilon_\mathrm{s}$}]
            \draw[hscink!60, line width=0.5pt] (axis cs:0,0) -- (axis cs:20,0);
            \addplot[hscseriesB, line width=1.2pt] coordinates
              {(0,-1) (5,-1) (5,0) (10,0) (10,1) (15,1) (15,0) (20,0)};
          \end{axis}
        \end{tikzpicture}\par}\smallskip
        (The emf could also be drawn positive first, then negative. What matters is that
        the two pulses have opposite signs and the same size.)\par
        (b) From \qty{5}{ms} to \qty{10}{ms} the current is steady, so the flux in the
        core is constant. Then $\dfrac{\Delta\Phi}{\Delta t} = 0$, so by Faraday's law the
        emf is zero. This is why a battery (steady DC) gives only a short pulse in the
        secondary when it is switched on or off.
        \begin{markingcriteria}
          \criterion{Sketch: constant emf from 0 to \qty{5}{ms}, zero from \qty{5}{ms} to
            \qty{10}{ms}, constant emf of the same size and opposite sign from
            \qty{10}{ms} to \qty{15}{ms}; explains the zero with a constant flux, so
            $\dfrac{\Delta\Phi}{\Delta t} = 0$ (Faraday's law)}{3}
          \criterion{Correct sketch, or correct explanation with one error in the
            sketch}{2}
          \criterion{One correct feature (for example, zero emf while the current is
            steady)}{1}
        \end{markingcriteria}
  \item $I_\mathrm{s} = I_\mathrm{p}\dfrac{N_\mathrm{p}}{N_\mathrm{s}}
        = \qty{50}{A} \times \dfrac{1}{25} = \qty{2.0}{A}$.
        Answer: \final{$I_\mathrm{s} = \qty{2.0}{A}$}. The voltage goes up 25 times, so
        the current goes down 25 times.
  \item A step-down transformer has fewer secondary turns, so $V_\mathrm{s}$ is small.
        Power in = power out ($V_\mathrm{p}I_\mathrm{p} = V_\mathrm{s}I_\mathrm{s}$), so a
        small $V_\mathrm{s}$ means a large $I_\mathrm{s}$. The secondary current is
        $\dfrac{N_\mathrm{p}}{N_\mathrm{s}}$ times the primary current. The large current
        heats and melts the metal, while the low voltage is safer for the welder.
        \begin{markingcriteria}
          \criterion{Fewer secondary turns give a small $V_\mathrm{s}$; power in = power
            out, so the secondary current is large; the large current heats and melts the
            metal}{3}
          \criterion{Two of these points, linked}{2}
          \criterion{States that a step-down transformer gives a larger current}{1}
        \end{markingcriteria}
  % ---- 4.2 Real transformers
  \item $\eta = \dfrac{P_\mathrm{out}}{P_\mathrm{in}} \times 100\%
        = \dfrac{470}{500} \times 100\% = 94\%$. Power wasted
        $= \qty{500}{W} - \qty{470}{W} = \qty{30}{W}$. In 10 minutes
        ($\Delta t = \qty{600}{s}$):
        $\Delta E = P\,\Delta t = \qty{30}{W} \times \qty{600}{s} = \qty{1.8e4}{J}$.
        Answers: \final{$\eta = 94\%$} and \final{$\Delta E = \qty{1.8e4}{J}$}.
  \item The iron core is a conductor, and the AC makes the flux in it change all the
        time. By Faraday's law this induces eddy currents that circulate inside the core.
        By Lenz's law, their magnetic field opposes the change in flux. They heat the core
        ($I^2R$), so energy is wasted. A laminated core is made of thin iron sheets with
        insulation between them. The sheets lie along the flux. The insulation confines
        each eddy current to one thin sheet. So each loop surrounds less of the changing
        flux (a smaller emf), and it has a narrow path with a high resistance. So the eddy
        currents and the heating are much smaller, and the transformer is more efficient.
        \begin{markingcriteria}
          \criterion{The changing flux induces eddy currents in the conducting core, which
            heat it; thin insulated sheets along the flux confine them to thin,
            high-resistance paths with a small emf, so the eddy currents and the heating
            are much smaller}{3}
          \criterion{One of these links explained (how eddy currents waste energy, or how
            laminations reduce them)}{2}
          \criterion{States that laminations reduce eddy currents}{1}
        \end{markingcriteria}
  \item The limitation is \final{incomplete flux linkage}. With no core, much of the
        flux made by the primary spreads out into the air and never passes through the
        secondary (\cref{ch4:fig:linkage}(a)). So the flux through each secondary turn
        changes less, and the induced emf is smaller than the turns equation predicts.
        To fix it: wrap both coils around one closed soft-iron core, or wind the
        secondary directly over the primary.
        \begin{markingcriteria}
          \criterion{Names incomplete flux linkage; explains that much of the primary's
            flux misses the secondary, so its emf is smaller; gives a correct fix}{3}
          \criterion{Two of these points}{2}
          \criterion{One of these points}{1}
        \end{markingcriteria}
  \item $P_\mathrm{in} = V_\mathrm{p}I_\mathrm{p} = \qty{230}{V} \times \qty{2.0}{A}
        = \qty{460}{W}$. $P_\mathrm{out} = \eta P_\mathrm{in} = 0.85 \times \qty{460}{W}
        = \qty{391}{W}$. $I_\mathrm{s} = \dfrac{P_\mathrm{out}}{V_\mathrm{s}}
        = \dfrac{\qty{391}{W}}{\qty{46}{V}} = \qty{8.5}{A}$.
        Answer: \final{$I_\mathrm{s} = \qty{8.5}{A}$} (the ideal value would be
        \qty{10}{A}).
  \item There is no shared closed iron core to guide the flux, and there is an air gap.
        So much of the changing flux from the pad's coil misses the phone's coil
        (incomplete flux linkage). So less energy is passed to the phone. To deliver
        the same power, the pad must drive a larger current, so more energy becomes
        $I^2R$ heat in the coils. The leaked changing flux can also induce eddy currents
        in nearby metal (parts of the phone, a coin), which heats it. So more energy is
        wasted as heat, and the pad and phone get warm. (Lining up the coils well, and
        ferrite sheets behind the coils to guide the flux, reduce the waste.)
        \begin{markingcriteria}
          \criterion{Air gap and no core, so much flux misses the phone's coil
            (incomplete flux linkage); a larger current is needed, so more $I^2R$ heating
            (or eddy-current heating in nearby metal)}{3}
          \criterion{Explains flux leakage, or links a larger current to more
            heating}{2}
          \criterion{States that more flux leaks}{1}
        \end{markingcriteria}
  % ---- 4.3 Sending electricity at high voltage
  \item $P_\mathrm{loss} = I^2R = (\qty{150}{A})^2 \times \qty{2.0}{\ohm}
        = \qty{45000}{W}$. Answer: \final{$P_\mathrm{loss} = \qty{4.5e4}{W}$}
        (\qty{45}{kW}).
  \item $P = VI$ stays the same, so 4 times the voltage gives $\tfrac{1}{4}$ of the
        current. The loss is $I^2R$, so it becomes $\left(\tfrac{1}{4}\right)^2
        = \tfrac{1}{16}$ of the old loss. Answer: \final{current $\times\tfrac14$, loss
        $\times\tfrac{1}{16}$}.
  \item The student used the line voltage (\qty{330}{kV}) in $\dfrac{V^2}{R}$. In that
        formula $V$ must be the voltage \emph{drop} across the line, not the voltage sent.
        The wrong answer, \qty{1.8e10}{W}, is about 90 times the power sent, which is
        impossible. Correct method:
        $I = \dfrac{P}{V} = \dfrac{\qty{2.00e8}{W}}{\qty{3.30e5}{V}} = \qty{606}{A}$, then
        $P_\mathrm{loss} = I^2R = (\qty{606}{A})^2 \times \qty{6.0}{\ohm}
        = \qty{2.20e6}{W}$. Answer: \final{$P_\mathrm{loss} = \qty{2.2}{MW}$} (1.1\% of
        the power sent). Check: the real voltage drop is
        $IR = \qty{606}{A} \times \qty{6.0}{\ohm} = \qty{3636}{V} \approx \qty{3.6}{kV}$,
        and $\dfrac{(\qty{3636}{V})^2}{\qty{6.0}{\ohm}} = \qty{2.2e6}{W}$, the same
        answer.
  \item Why AC: transformers work by electromagnetic induction, which needs a changing
        flux. AC makes a flux that changes all the time, so its voltage can be stepped up
        for the long lines and down again. Steady DC gives a constant flux, so a
        transformer cannot change its voltage. Why step up: for a fixed power,
        $P = VI$. So a higher voltage means a smaller current. The loss is $I^2R$, so it
        falls with the square of the current. Why step down: very high voltages are
        dangerous. They can arc across air and need big insulators. Also, appliances are
        built for \qty{230}{V}.
        \begin{markingcriteria}
          \criterion{Transformers need AC (a changing flux) to step the voltage up and
            down}{1}
          \criterion{For a fixed power, $P = VI$, so a higher voltage gives a smaller
            current}{1}
          \criterion{$P_\mathrm{loss} = I^2R$, so the loss falls with the square of the
            current}{1}
          \criterion{The voltage is stepped down for safety and because appliances need
            \qty{230}{V}}{1}
        \end{markingcriteria}
  \item (a) Voltage drop along the line:
        $V_\mathrm{drop} = IR = \qty{200}{A} \times \qty{5.0}{\ohm} = \qty{1000}{V}
        = \qty{1.0}{kV}$. Voltage at the far end:
        $\qty{67.0}{kV} - \qty{1.0}{kV} = \qty{66.0}{kV}$.
        Answer: \final{\qty{66.0}{kV}}.\par
        (b) The output (\qty{11.0}{kV}) is smaller than the input (\qty{66.0}{kV}), so it
        is a \final{step-down} transformer.
        $\dfrac{N_\mathrm{p}}{N_\mathrm{s}} = \dfrac{V_\mathrm{p}}{V_\mathrm{s}}
        = \dfrac{\qty{66.0}{kV}}{\qty{11.0}{kV}} = 6$, so \final{$N_\mathrm{p} :
        N_\mathrm{s} = 6 : 1$}.\par
        (c) The power that reaches the transformer is
        $P = VI = \qty{6.60e4}{V} \times \qty{200}{A} = \qty{1.32e7}{W}$. Ideal, so the
        same power comes out:
        $I_\mathrm{s} = \dfrac{P}{V_\mathrm{s}} = \dfrac{\qty{1.32e7}{W}}{\qty{1.10e4}{V}}
        = \qty{1200}{A}$. Answer: \final{$I_\mathrm{s} = \qty{1.20e3}{A}$}. Check:
        $I_\mathrm{s} = I_\mathrm{p}\dfrac{N_\mathrm{p}}{N_\mathrm{s}}
        = \qty{200}{A} \times 6 = \qty{1200}{A}$ \cmark. The current goes up because the
        voltage goes down.
\end{chapteranswers}
